% Generated draft: requires mathematical and source review.
\documentclass[11pt]{article}
\usepackage[margin=25mm]{geometry}
\usepackage{fontspec}
\setmainfont{Latin Modern Roman}
\usepackage{amsmath,amssymb,mathtools,unicode-math}
\setmathfont{Latin Modern Math}
\usepackage{enumitem,needspace,xurl,hyperref}
\setlength{\parindent}{0pt}
\setlength{\parskip}{5pt}
\setlength{\emergencystretch}{3em}
\newcommand{\N}{\mathbb{N}}
\newcommand{\Z}{\mathbb{Z}}
\newcommand{\Q}{\mathbb{Q}}
\newcommand{\R}{\mathbb{R}}
\newcommand{\C}{\mathbb{C}}
\newcommand{\F}{\mathbb{F}}
\newcommand{\A}{\mathbb{A}}
\newcommand{\PP}{\mathbb{P}}
\newcommand{\E}{\mathbb{E}}
\newcommand{\eps}{\varepsilon}
\newcommand{\dd}{\,\mathrm{d}}
\begin{document}
\section*{UnsolvedMath ID: 30000966}
\textbf{Source problem code:} catalogue entry
\subsection{Definitions and mathematical statement}
A stationary random point field on \(\R^d\) can be represented by a random locally finite counting measure \(\Xi\) whose law is invariant under translations. For a bounded Borel set \(A\), write \(N(A)=\Xi(A)\) for its number of points. When the intensity \(\rho\) is finite, stationarity gives the usual centering \(N(A)-\rho|A|\); finite intensity is not explicitly assumed in the supplied record.

The report-level question is: what large-deviation behavior governs bulk fluctuations of stationary random point fields, as distinguished from the boundary-effect fluctuations described in the source for the number of points in a large disk? A large-deviation formulation requires specifying the point-field law, the bulk observable or event, its scaling, a speed tending to infinity, and a rate function or an asymptotic probability. None of these is fixed by the extraction. In particular, stationarity alone is not a specification of one universal bulk law, and the disk-count comparison does not identify a particular process or precise boundary-fluctuation theorem. This remains an exploratory question; no particular hole-probability asymptotic, counting-function principle, or beta-ensemble model is silently substituted.
\subsection{Short English statement}
Describe large deviations of bulk fluctuations in the intended stationary point fields, beyond the boundary-driven disk-count fluctuations mentioned by the source. The process, bulk event, and asymptotic scale need to be recovered before a single precise large-deviation theorem can be stated.
\subsection{Sources}
\begin{itemize}
\item Local UnsolvedMath registry, permanent ID \(30000966\), source code \texttt{catalogue entry}, ``Bulk Large Deviations in Stationary Point Fields.'' Registry curation: Ulam AI / UnsolvedMath Contributors, CC BY 4.0; underlying sources retain their supplied provenance.
\item \emph{Stochastic Analysis} (2009), Oberwolfach Reports, pp.\ 1361--1424, DOI \url{https://doi.org/10.4171/owr/2008/25}.
\item Supplied ``Valko--Virag large gaps theorem'': \url{https://arxiv.org/abs/0811.0007}.
\item The supplied literature triage checked 2026-08-21 labels the problem partially solved: bulk large-gap and related large-deviation asymptotics for beta-ensemble limits have been proved, while the broader report-level bulk question remains. This is a dated source assessment, not an independent status check.
\end{itemize}
\end{document}
